\section{Do Dogs Go to Heaven}

\paragraph{Like a Moth to a Flame}

The two primary conditions of elementary life are attraction and repulsion. A healthy organism will be attracted to whatever is good for it, including its genetic integrity. Likewise, it will be repulsed by anything that threatens its well-being, again including its genetic integrity.

In the animal soul, an emotional element is added to those conditions. Attraction and repulsion are no longer merely mechanical, but are reinforced by that strong emotional component. Since there is no abstract intelligence in animals, these emotions can be easily fooled. For example, the proverbial moth is attracted to light, which is good, but also to a flame, which will be its downfall.

Many people live primarily an animal life of emotion, with occasional forays into a true intellectual understanding. These emotions have quite a hold, so it is difficult to dissuade them from unnatural attractions and repulsions, because they feel so good. That emotional life is so sweet and appealing, except when it isn't.

The definition of love, for such people, is defective, since it is purely passive:

\begin{quotex}
Love is attraction with a strong and pleasant emotional component.

\end{quotex}
True Love, on the other hand, is active, a matter of the will. The intellect must decide what is attractive and repulsive; the emotions and body will follow.

\begin{quotex}
To Love is to will the Good.

\end{quotex}
\paragraph{Stray Dogs in the Intellect}
\begin{quotex}
If our intellect is inexperienced in the art of watchfulness it at once begins to entertain whatever impassioned fantasy appears in it, and plies it with illicit questions and responds to it illicitly. Then our own thoughts are conjoined to the demonic fantasy, which waxes and burgeons until it appears lovely and delectable to the welcoming and despoiled intellect. The intellect then is deceived in much the same way as lambs when a stray dog comes into the field in which they happen to be: \emph{in their innocence they often run towards the dog as though it were their mother, and their only profit in coming near it is that they pick up something of its stench and foulness}.

In the same way our thoughts run ignorantly after demonic fantasies that appear in our intellect and, as I said, the two join together and one can see them plotting to destroy the city of Troy like Agamemnon and Menelaus. For they plot together the course of action they must take in order to bring about, in practice and by means of the body, that purpose which the demons have persuaded them is sweet and delectable. In this Way sins are produced in the soul: and hence the need to bring out into the open what is in our hearts. \flright{\textsc{Philokalia}}

\end{quotex}
\paragraph{Beyond Good and Evil}
Many people claim to have learned of “unconditional love” from their relationship with a pet. Now unconditional love is not necessarily desirable if the object of such love is not the Good. True Love is an act of the will toward the Good.

Now the dog is beyond good and evil, and has no interest in the moral state of its human companion. It is only interested in what is good or bad for its own well-being.

Oddly enough, such people like being the recipient of the pet's “unconditional love”, but not so much as the giver of such love. A woman who took that lesson to heart would experience better relations with her lover. For example, she would give him three square meals a day, clean up his shit without complaint, and scratch his belly every night. Try that at home and let me know.

\paragraph{States of the Spirit in the World}
In \emph{Heaven and Hell} \#435, \textbf{Emmanuel Swedenborg} writes:

\begin{quotex}
All this has been said to convince the rational man that viewed in himself man is a spirit, and that the corporeal part that is added to the spirit to enable it to perform its functions in the natural and material world is not the man, but only an instrument of his spirit. But evidences from experience are preferable, because there are many that fail to comprehend rational deductions; and those that have established themselves in the opposite view turn such deductions into grounds of doubt by means of reasonings from the fallacies of the senses. 

\end{quotex}
Here is the key to understand Swedenborg. He is not really describing a visit to the spirits in Heaven or Hell at some future time, but he is describing phenomenologically the spiritual states of people in the here and now. This is a matter of direct experience and is called in Tradition “external considering” or “reading souls”.

This is obvious from this paragraph from the web site of the Swedenborg Foundation:

\begin{quotex}
Swedenborg writes that each person self-selects a destination in the afterlife depending on what they loved most while on earth. Those who become evil spirits are people who love things like power and wealth and even being cruel to others at the cost of all else. After death, such people push away anything that was good within them and become completely focused on the thing that they value above all else—they become their emotions in tangible form. Spirits are grouped together into societies according to what they love.

\end{quotex}
Such behavior does not just happen “after death”, but right now in the life of such persons. The death referred to here is a “spiritual death”. Swedenborg's colorful descriptions of the various evil spirits he encountered are his way of describing their inner soul life.

Continuing his discussion of \#435, Swedenborg denies that pets “go to heaven” because, unlike man, they are not spirits. Nevertheless, many people have little self-knowledge of themselves as spirits in the world and remain entrenched in corporeal imagery. They can see little difference between themselves and animals.

\begin{quotex}
People who have convinced themselves [of the opposite view] tend to think that animals live and sense just the way we do, so that they too have a spiritual nature like ours; yet this dies along with their bodies. However, the spiritual nature of animals is not the same as ours. We have an inmost nature that animals do not, a nature into which the Divine flows and which it raises toward itself, in this way uniting us to itself. So we, unlike animals, can think about God and about divine matters of heaven and the church. We can love God because of these matters and by engaging with them; and so be united to him; and anything that can be united to the Divine cannot be destroyed. Anything that cannot be united to the Divine, though, does disintegrate.”

\end{quotex}
I don't believe he writes about the destination of those who love their dogs more than they love God, other than that they will organize themselves together, probably into a pack.

\paragraph{Living like a Dog}
A year or two ago, I watched a short feature — whose title now eludes me — about a man who designed a pharmaceutical that turned him temporarily into a dog. His drug became very popular, so that clubs arose whose members became dogs, prowling the streets of the city, and rutting, for a night, only to return to the human state in the morning.

Eventually, the authorities regarded his drug as the cause of a public nuisance and ordered him to cease and desist manufacturing and distributing the drug. Instead, he found a way to make the effects permanent and turned himself — and his many followers — into dogs.

\paragraph{The Complete Man}
\textbf{Nicolas Cabasilas} (1320-1392) is a Greek saint and his writings appear in the Roman Liturgy of the Hours. According to \textbf{Mircea Eliade}, he considered the layman to be superior to the monk. The goal of the latter is the \emph{angelic life} while that of the layman is the \emph{complete man}.

\begin{quotex}
[The Law of Love] demands no arduous nor afflicting work, nor loss of money; it does not involve shame, nor any dishonor, nor anything worse; it puts no obstacle in the pursuit of any art or profession. The general keeps the power to command, the labourer can work the ground, the artisan can carry on with his occupation. There is no reason to retire into solitude, to eat unusual food, to be inadequately clothed, or endanger one's health, or to resort to any other special endeavor; it suffices to give oneself wholly to meditation and to remain always within oneself without depriving the world of one's talents. \flright{\textsc{Nicolas Cabasilas}}

\end{quotex}
This confirms our claim about the need, in our time, for the spiritual man to be active in the world. The complete man actualizes all his possibilities, not just the angelic ones. The spirituality of the monks has dominated religious discourse, although it is not necessarily appropriate for the man living in the world. It is to these latter men that the future belongs.

\paragraph{Christendom and Skillful Means}
\textbf{Frithjof Schuon} describes the spiritual state of European man, and why he had to overcome the ancient pagan religion:

\begin{quotex}
European humanity has about it something Promethean and tragic. It therefore required a religion which could surpass and sublimate the dramatism of the Greek and Germanic gods and heroes. Furthermore, the European creative genius implies a need to “burn what one has worshipped,” and from this comes a prodigious propensity for denial and for change. The Renaissance offers us the clearest proof and the most stupefying example of this, not to speak of what is going on in our own time on an incomparably more dangerous level; it is still “Man” that is at issue, but with totally different emphases.

\end{quotex}
The new religion overcame the arbitrary forces, experienced as the whims of gods and goddesses, that seemed to dominate life. It added Providence to the forces of Destiny and Will, bringing warmth to the world. With the Logos, the world could be experienced as an ordered cosmos rather than the result of Fate or blind Will. Nevertheless, Schuon points to a certain tension, which has only become more exacerbated over time.

\begin{quotex}
Christianity has a dramatic quality about it: it has the sense of the Sublime rather than that of the Absolute, and the sense of Sacrifice rather than that of Equilibrium. In this second aspect it extends to society as a whole a vocation that is properly speaking ascetic – above all in the Latin Church – something which, as a particular \emph{upaya}, it certainly has a right to do, but which has nonetheless provoked historical disequilibriums that are both disastrous and providential.

\end{quotex}
Time, Space, and Equilibrium are the three preliminary conditions of manifestation from which the active, passive, and neutralizing forces arise. If Tradition is the Civilization of Space, then Time will tend to consume Space unless Equilibrium is established. The re-establishment of equilibrium will require the replacement of “skillful means”, that may have been appropriate at a given time and place, but are no longer effective. The tendency to accept upaya as equivalent to dogma must be avoided, since skillful means are appropriate to a particular time and place. Ultimately the goal is \emph{prajna}, i.e., gnosis.

\paragraph{The Eucharist and Nonduality}
\begin{quotex}
The Tao acts and does not exhaust itself. \flright{\textsc{Tao Te Ching}}

\end{quotex}
About 6 weeks ago, I heard a sermon at mass on the Eucharist by a visiting priest. In retrospect, I wish I had taken notes for the 45 minute talk, but I had no expectation of it.

Ultimately, it was a nondual explanation of the Eucharist. The thought experiment for the priest was how the Christ can be present in all the hosts while remaining undivided. He compared Christ to the Sun, which is inexhaustible, while all the masses received the light and substance of the Sun. However, this is not a dualism. The priest himself represents Christ, so that he is also present in the Eucharist. Moreover, the whole congregation, is also the Body of Christ, and so is present in the Eucharist.

The sacrifice of the Mass is not a repetition of Calvary, as the heretics mistakenly assert, but is rather the very \emph{same} sacrifice. So the Eucharist transcends both time and space, and is a theophany connecting the world of being to that of becoming.

\paragraph{The Conversion of Russia}
At the end of the Latin rite mass, prayers are said for the conversion of Russia. The missal explains:

\begin{quotex}
These prayers were introduced by Pope Leo XIII to obtain an acceptable solution to the Vatican's relations with the Italian State after the seizure of the Papal States. After its resolution by the establishment of the Vatican State through the Treaty of 1929, Pope Pius XI asked that these prayers should be said for the conversion of Russia.

\end{quotex}
So for over 80 years, these prayers have been said. Perhaps some have noticed a small change in Russia today, different from what it was in 1929. What Western ruler has asked his subordinates to study the works of such eminent philosophers like Vladimir Solovyov and Nicolas Berdyaev?

If the Orthodox still have a beef with the West, perhaps they should return the favor instead of engaging in useless carping.

\paragraph{Huxley on Love}
\textbf{Thomas Molnar}, in \emph{The Decline of the Intellectual}, notes the fundamental message of \textbf{Aldous Huxley}'s dystopian novel \emph{Brave New World}:

\begin{quotex}
Aldous Huxley had come to the disturbing conclusion that culture, that is, the inner man, cannot survive without tragedy, that is, outside traditional society, where human rapports and not scientific organization, prevail. This is the disillusioned statement of the World Comptroller from which everything incalculable, the last ounce of freedom and hence the restlessness accompanying creative work, had been eliminated.

\end{quotex}
That is, the World Comptroller wants knowledge without gnosis, by which he will organize society. So-called scientific control of society is the dream of leftism, although it is a nightmare for those controlled. It tries to control the outward circumstance of life to replace the sense of inwardness.

It leaves out of the picture transcendence, darkness, the depths — or the abyss. Consider, for example, \textbf{Jacob Boehme}'s revelation of the Ungrund which is prior to Being, from which creative freedom arises, which cannot be anticipated in any scientific organization. The world is a mix, then, of Darkness and Light, and the darkness cannot be hidden by any technocrat. Nay, he actually is the darkness.

Unfortunately, Huxley eventually ended up California dreaming, that strange attractor. Molnar explains:

\begin{quotex}
In a series of lectures sponsored by the University of California, he advocated the strengthening of love among people by suggesting that mothers all over the world, while suckling their babes should rub them against other people and animals saying “nice-nice, good-good,” and should point at everybody around repeating some incantatory formula like “this is a good man, we must love him.”

\end{quotex}
In other words, he thinks people should start acting more like pet dogs.

\paragraph{Loving your Enemies}
If true love is to will the good for someone, then the desire for justice, as the greatest natural virtue, is an expression of love. For example, a general would engage only in a just war with an enemy, not targeting women and children, nor anything without a military purpose. Such a general would become a complete man.

Now the principle of justice is “to each his own”, or in other words, in justice a man should get what he is entitled to. In particular, a good general will give the enemy exactly what he deserves.

\paragraph{Answer at the Back of the Book}
\begin{quotex}
The angels are the thoughts of God incarnated in the celestial flesh. \flright{\textsc{Pierre Deghaye}, \emph{La Naissance de Dieu}}

\end{quotex}
Specifically, there is an angel of the dogs, which serves as the “group I” for all dogs, which lack an individual “I”. This angel will know all the dogs, in a supra-individual way, as non-formal manifestation. Presumably you can communicate with that angel and experience your pet “from the inside”, in a much more intimate way.

But first of all, you need to work out your own salvation.



\flrightit{Posted on 2016-09-07 by Cologero }

\begin{center}* * *\end{center}

\begin{footnotesize}\begin{sffamily}



\texttt{William on 2016-09-08 at 00:33 said: }

“According to Mircea Eliade, he considered the layman to be superior to the monk.”

Does Gornahoor also take such a view? I understand the need to actualize possibilities beyond the angelic, but how can the assertion of lay superiority over the monastics be reconciled with the eminence of the priestly caste, which is clearly above the laymen in the caste system.


\hfill

\texttt{Cologero on 2016-09-08 at 07:58 said: }

So are you “off the hook”, William? Do you want a free ride, while letting the monks do all that hard stuff?

We accept them as teachers, so we posted a long excerpt from the Philokalia. How will you use that passage in your life today?

By the way, not all monks are priests.

In a very recent post, we made this suggestion in regard to some ideas from Joachim di Fiore:

\begin{quotex}
Hence, the Age of the Holy Spirit is actually the beginning of a new cycle, not the final stage of some evolutionary process. If the first age was that of the fathers, then of priests, Joachim thought the Third Age would be that of an elite corps of contemplative monks. Actually, it is more likely that that corps will consist of men and women active in the world, while transcending it. 

\end{quotex}
\texttt{William on 2016-09-08 at 19:46 said: }

My question has nothing to do with whether or not I'm off the hook. I know what my duties are in accordance with the caste system. I'm simply asking how we can consider the priestly caste to be higher in the hierarchy than most of us will ever reach, yet still consider ourselves as lay people to be superior to monastic priests. Also, what are we to make of lay brothers? Perhaps they are who the sentence in question was referring to and I'm just misunderstanding? Even then, though, wouldn't they be considered ascetics, which Evola described as a path of the Absolute Man?


\hfill

\texttt{Pepé on 2016-09-08 at 20:59 said: }

William – as I understand it, that means a `new caste'. so to speak.


\hfill

\texttt{Cologero on 2016-09-09 at 00:01 said: }

The castes are all mixed up today, and few may be in their correct caste. So your point is irrelevant.


\hfill

\texttt{mleow on 2016-09-09 at 03:06 said: }

Moving away from christianity, sanatana dharma which unarguably has more authority to it (from a traditional perspective) then some sayings from swedenborg, considers that animals definitely do have a soul life for example a short summary of the views expressed;

http://www.hinduwebsite.com/hinduism/essays/animals.asp

Besides cologero it is almost amusingly predictable how your canine hatred shines trough and resurface in every odd post you made here over the years. Were you bitten by a dog when you were young or something? Your knowledge and experience of dogs seems to be limited, “Now the dog is beyond good and evil, and has no interest in the moral state of its human companion. It is only interested in what is good or bad for its own well-being.” Is it really cats you are referring?? St Roch would probably not concur with that statement..


\hfill

\texttt{Cologero on 2016-09-09 at 08:24 said: }

No, mleow, but I did bite a dog once.


\hfill

\texttt{????? on 2016-09-09 at 10:13 said: }

I recall hearing someone ask a similar or the same question at a sufi gathering, and the response concerned the dog who lived in the cave with the (seven?) sleepers. Since the dog was a good dog and a companion of saints, it will live on as a human in Paradise insha'Allah. IIRC. Because you will be with those whom you love. But will dogs be with those whom they love? That hadith wasn't directly addressed to dogs… This is a controversial issue among Muslims too.

\texttt{It was narrated that Abu Hurayrah (may Allah be pleased with him) said: The Messenger of Allah (peace and blessings of Allah be upon him) said: “Pray in the sheep pens and wipe their dust (raghaam), for they are among the animals of Paradise.”}


\hfill

\texttt{Ulysses on 2016-09-09 at 13:44 said: }

On the subject of experiencing pets “from the inside”, Groucho Marx said “Outside of a dog, a book is man's best friend. Inside of a dog, it's too dark to read.”


\hfill


\hfill


\hfill

\texttt{Matt on 2016-09-09 at 23:11 said: }

Mleow,

The Church also recognizes that animals have a soul. The question is, do they have an intellectual/rational soul? Because the matter of immortality rests with the intellectual soul. From Tradition (chritstian or not), humans have all three soul layers: vegitative, animal/sensitive, and intellectual/rational. That's what makes us fundamentally different from animals and plants (a difference of kind, and not merely of degree). Swedenborg's statements are in keeping with the authoritative teachings of the Church, and also with those of the other legitimate traditions.

Just putting aside Tradition and doctrine, and instead looking at empircal studies (if that's what matters to some readers), the supposed evidence for certain animals possessing substantive(moral) reason and intelligence is rather limited and weak, with the concepts of reason and intelligence being poorly grasped by the researchers that make claims about certain animals.

Ultimately , how does one envision what eternal life is like? Is it the ever deepening-ascending contemplation and love of the Divine, or is it forever playing fetch with Fido?


\hfill

\texttt{William on 2016-09-10 at 11:27 said: }

“The castes are all mixed up today, and few may be in their correct caste. So your point is irrelevant.”

I've been reading Gornahoor for quite a while, and out of all the asinine replies you've given to readers, this has to be the least satisfactory one yet.


\hfill

\texttt{Cologero on 2016-09-10 at 14:39 said: }

Apparently, William, my source is equally asinine:

\begin{quotex}
Have we not arrived at that terrible age, announced in the Sacred Books of India, `when the castes shall be mingled, when even the family shall no longer exist'.

As we have already pointed out, under the present state of affairs in the Western world, \textbf{nobody any longer occupies the place that he should normally occupy by virtue of his own nature; this is what is meant by saying that the castes no longer exist, for caste, in its traditional meaning, is nothing other than individual nature}, with the whole array of special aptitudes that this carries with it and that predisposes each man to the fulfillment of one or another particular function. 

\end{quotex}
From \emph{Crisis of the Modern World} by \textbf{Rene Guenon}


\hfill

\texttt{AC on 2016-09-10 at 18:48 said: }

Two relevant excerpts:

According to the [evolutionist explanation of the instinct of animals], instinct is the expression of the heredity of a species, of an accumulation of analogous experiences down the ages. This is how they explain, for example, the fact that a flock of sheep hastily gathers together around the lambs the moment it perceives the shadow of a bird of prey, or that a kitten while playing already employs all the tricks of a hunter, or that birds know how to build their nests. In fact, it is enough to watch animals to see that their instinct has nothing of an automatism about it. The formation of such a mechanism by a purely cumulative . . . process is highly improbable, to say the least. Instinct is a nonreflective modality of the intelligence; it is determined, not by a series of automatic reflexes, but by the “form”—the qualitative determination—of the species. This form is like a filter through which the universal intelligence is manifested. . . The same is also true for man: his intelligence too is determined by the subtle form of his species. This form, however, includes the reflective faculty, which allows of a singularization of the individual such as does not exist among the animals. Man alone is able to objectivize himself. He can say: “I am this or that.” He alone possesses this two-edged faculty. Man, by virtue of his own central position in the cosmos, is able to transcend his specific norm; he can also betray it, and sink lower; “The corruption of the best is corruption at its worst.” A normal animal remains true to the form and genius of its species; if its intelligence is not reflective and objectifying, but in some sort existential, it is nonetheless spontaneous; it is assuredly a form of the universal intelligence even if it is not recognized as such by men who, from prejudice or ignorance, identify intelligence with discursive thought exclusively. 

{}-Titus Burckhardt, Modern Psychology

And:

We have seen reason to believe that not all animals suffer as we think they do: but some, at least, look as if they had selves, and what shall be done for these innocents? 

The real difficulty about supposing most animals to be immortal is that immortality has almost no meaning for a creature which is not “conscious” in the sense explained above. If the life of a newt is merely a succession of sensations, what should we mean by saying that God may recall to life the newt that died to-day? It would not recognise itself as the same newt; the pleasant sensations of any other newt that lived after its death would be just as much, or just as little, a recompense for its earthly sufferings (if any) as those of its resurrected – I was going to say “self”, but the whole point is that the newt probably has no self. The thing we have to try to say, on this hypothesis, will not even be said. There is, therefore, I take it, no question of immortality for creatures that are merely sentient. Nor do justice and mercy demand that there should be, for such creatures have no painful experience. Their nervous system delivers all the letters A, P, N, I, but since they cannot read they never build it up into the word PAIN. And all animals may be in that condition. If, nevertheless, the strong conviction which we have of a real, though doubtless rudimentary, selfhood in the higher animals, and specially in those we tame, is not an illusion, their destiny demands a somewhat deeper consideration. The error we must avoid is that of considering them in themselves. Man is to be understood only in his relation to God. The beasts are to be understood only in their relation to man and, through man, to God. Let us here guard against one of those untransmuted lumps of atheistical thought which often survive in the minds of modern believers. Atheists naturally regard the co-existence of man and the other animals as a mere contingent result of interacting biological facts; and the taming of an animal by a man as a purely arbitrary interference of one species with another. The “real” or “natural” animal to them is the wild one, and the tame animal is an artificial or unnatural thing. But a [believer] must not think so. Man was appointed by God to have dominion over the beasts, and everything a man does to an animal is either a lawful exercise, or a sacrilegious abuse, of an authority by divine right. The tame animal is therefore, in the deepest sense, the only “natural” animal – the only one we see occupying the place it was made to occupy, and it is on the tame animal that we must base all our doctrine of beasts. Now it will be seen that, in so far as the tame animal has a real self or personality, it owes this almost entirely to its master. If a good sheepdog seems “almost human” that is because a good shepherd has made it so. . . I am now going to suggest that . . . beasts that attain a real self are \_in\_ their masters. That is to say, you must not think of a beast by itself, and call that a personality and then inquire whether God will raise and bless that. You must take the whole context in which the beast acquires its selfhood – namely “The good man – and – the – goodwife – ruling – their children – and – their – beasts – in – the – good – homestead”. That whole context may be regarded as a “body”; and how much of that “body” may be raised along with the goodman and the goodwife, who can predict? So much, presumably, as is necessary not only for the glory of God and the beatitude of the human pair, but for that particular glory and that particular beatitude which is eternally coloured by that particular terrestrial experience. And in this way it seems to me possible that certain animals may have an immortality, not in themselves, but in the immortality of their masters. And the difficulty about personal identity in a creature barely personal disappears when the creature is thus kept in its proper context. If you ask, concerning an animal thus raised as a member of the whole Body of the homestead, where its personal identity resides, I answer “Where its identity always did reside even in the earthly life – in its relation to the Body and, specially, to the master who is the head of that Body”. In other words, the man will know his dog, the dog will know its master and, in knowing him, will be itself. To ask that it should, in any other way, know itself, is probably to ask for what has no meaning. Animals aren't like that, and don't want to be. 

My picture of the good sheepdog in the good homestead does not, of course, cover wild animals nor (a matter even more urgent) illtreated domestic animals. But it is intended only as an illustration drawn from one privileged instance – which is, also, on my view the only normal and unperverted instance of the general principles to be observed in framing a theory of animal resurrection. I think [believers] may justly hesitate to suppose any beasts immortal; for two reasons. Firstly, because they fear, by attributing to beasts a “soul” in the full sense, to obscure that difference between beast and man which is as sharp in the spiritual dimension as it is hazy and problematical in the biological. And secondly, a future happiness connected with the beast's present life simply as a compensation for suffering – so many millenniums in the happy pastures paid down as “damages” for so many years of pulling carts, seems a clumsy assertion of Divine goodness. We, because we are fallible, often hurt a child or an animal unintentionally, and then the best we can do is to “make up for it” by some caress or titbit. But it is hardly pious to imagine omniscience acting in that way – as though God trod on the animals' tails in the dark and then did the best He could about it! In such a botched adjustment I cannot recognise the master-touch; whatever the answer is, it must be something better than that. The theory I am suggesting tries to avoid both objections. It makes God the centre of the universe and man the subordinate centre of terrestrial nature: the beasts are not co-ordinate with man, but subordinate to him, and their destiny is through and through related to his. And the derivative immortality suggested for them is not a mere amende or compensation: it is part and parcel of the new heaven and new earth, organically related to the whole suffering process of the world's fall and redemption. Supposing, as I do, that the personality of the tame animals is largely the gift of man – that their mere sentience is reborn to soulhood in us as our mere soulhood is reborn to [God] – I naturally suppose that very few animals indeed, in their wild state, attain to a “self” or ego. But if any do, and if it is agreeable to the goodness of God that they should live again, their immortality would also be related to man – not, this time, to individual masters; but to humanity. That is to say, if in any instance the quasi-spiritual and emotional value which human tradition attributes to a beast (such as the “innocence” of the lamb or the heraldic royalty of the lion) has a real ground in the beast's nature; and is not merely arbitrary or accidental, then it is in that capacity, or principally in that, that the beast may be expected to attend on risen man and make part of his “train”. 

When we are speaking of creatures so remote from us as wild beasts, and prehistoric beasts, we hardly know what we are talking about. It may well be that they have no selves and no sufferings. It may even be that each species has a corporate self–that Lionhood, not lions, has shared in the travail of creation and will enter into the restoration of all things. And if we cannot imagine even our own eternal life, much less can we imagine the life the beasts may have as our “members”. If the earthly lion could read the prophecy of that day when he shall eat hay like an ox, he would regard it as a description not of heaven, but of hell. And if there is nothing in the lion but carnivorous sentience, then he is unconscious and his “survival” would have no meaning. But if there is a rudimentary Leonine self, to that also God can give a “body” as it pleases Him–a body no longer living by the destruction of the lamb, yet richly Leonine in the sense that it also expresses whatever energy and splendour and exulting power dwelled within the visible lion on this earth. I think, under correction, that the prophet used an eastern hyperbole when he spoke of the lion and the lamb lying down together. That would be rather impertinent of the lamb. To have lions and lambs that so consorted would be the same as having neither lambs nor lions.”

{}-CS Lewis, The Problem of Pain


\hfill

\texttt{Max on 2016-09-11 at 05:55 said: }

We have anticipated the discussion of this intriguing question since it first appeared among the upcoming topics.

There are “mysteries” of the human condition that demand an answer as well, for example why some people are reduced to speaking baby talk as soon as they see a dog. I heard children of parents who speak to them in such a way grow up to become less intelligent and capable, so one can legitimately ask who is really benefiting from it since it is obviously not the children or the dogs? People may try to convince themselves that it is a benevolent concession to others, but it is more likely to correspond to their own projected needs.

Dogs does not understand “economy” and will not obey for moral reasons or to gain an advantage since they even stays with an owner who abuses them. Humans on the other hand, understanding economy, are always looking to gain one thing or another. What that is and how successful we are at it differs. I do not think that a hypothetical “dog heaven” would be very appealing, but perhaps that is what some people aspire to.

Dogs do not choose whether to “love” or not, it is an instinct that they need for survival since they are not even capable of hunting their own meals. Humans are also generally a domesticated race that does not need to hunt for survival because of setting in place welfare systems providing for them. What a dog feels for a human provider is like humans “unconditional” loyalty to these institutions. This tamed nature gets expressed for example as childlike acting or being drawn toward increased socialization, the seeking of approval from the environment, or speaking in baby talk which is the equivalent of a dog's immature bark retained throughout life (a mature wolf will not bark, it launches its deadly attack without warning, something that stupid zookeepers have repeatedly got the opportunity to experience).

The process as a whole involves certain advantages but it causes a higher dependency on the right upbringing, training, and so on to result in a well equipped and highly functional individual. Basically it allows for more “plasticity” and through that potential for freedom and self-direction, but at the same time also the risk of slavery, making it a kind of gamble. For example, the mental resources that are normally used for social approval can be redirected to the pursuit of knowledge. That does however not imply a lack of “social intelligence”, but rather the judgement of the direction (or lack thereof) of the majority as a meaningless evolutionary dead end. Humans mature at varying rates, and the longer the stage of childhood, the higher the potential intelligence of the mature man becomes, but with it the vulnerable formation period increases.

In North America, genetic testing has shown the existence of wild wolves bred with coyotes and domesticated dogs. It was speculated that this has resulted in a far more dangerous creature, perhaps inheriting the wolfs capabilities and merciless killer instinct, now combined with the other parts disrespect and fearlessness for the weak modern man, sort of like the animal equivalent of stealing the Promethean fire, serving to introduce new instincts into the wild population. As a veritable “age of the wolf”, this would be like the Caucasian race loosing its universal altruistic side, bred out of existence for lack of bringing any advantage in the current climate. The popular narrative always risks becoming a self-fulfilling prophecy when its inherent destiny proceeds to travel its course.

The technocratic attempt to minutely control the world has deprived our contemporaries of the feeling of a vast cosmos made up of various forces, wills, beings, an experience without which Christianity becomes difficult to approach. It seems that generally, and without an inevitable mechanical backlash to current trends bringing a deepened disclosure of reality, that higher view will remain closed to most.

Taxonomic species are often defined on the basis of whether they breed with each other or not. That however not only depends on physical opportunity but on will and intent, which could be both good or bad news depending on how you interpret it. For leftist egalitarians it would become possible to promiscuously become an “equal” to lower forms of life, which in turn means that we can without remorse exclude them from the species. Humanity seems to be on the brink not so much of “breakaway civilization” as breakaway speciation, that is to say shifts in Intelligence. We are seeing not the widespread surpassing of the totemic level (from “Social Facts and Group Control”), but a return to it. Since the modern welfare state is not based on an equilibrated state of “social fact” that can persist for centuries, perhaps in a hundred years, all that will be left of it is a primitive dog cult in the outback mindlessly chanting the Law of the Jungle in baby talk…

“And the Wolf that shall keep it may prosper, but the Wolf that shall break it must die.”

The article on epigenetics that you linked to (I cannot find on which page) raises the point that genes, or the entire genesis of an organism for that matter, is much more fluid than previously thought – it can quite literally absorb from its environment the genes that are congenial to its aspirations, even by other means than reproductive sex. All of this is consistent with the ancients account of for example “genii” or “daemon”, and also makes much more sense in a worldview of mind over matter.

It also makes the point that “survival of the fittest” merely means “survival of those who survive” which is circular reasoning with no explanatory power. Now, we may in an exercise of conscious judgement choose to measure “fitness” by other criteria than reproduction, acquiring a view not merely of a vague “survival”, but what lives on and in which way.

Epigenetics tells us that the expression of our genes, influencing what we become, is partly dependent on our environment. As a spontaneous experience, thought belongs to the environment of the human organism. This means in practice, even in a materialistic framework where thought is nothing but chemical reactions (those reactions are no less fundamental an environment than chemical reactions occurring outside the body), that thinking something specific holds the potential of unlocking and activating, through providing the correct stimuli, aspects of our inheritance that were previously dark and inaccessible.


\hfill

\texttt{Cologero on 2016-09-11 at 11:22 said: }

Thanks for finding those excerpts, AC.

The one from Burckhardt makes more sense, as it is closer to our viewpoint. The vegetable, animal, and intellectual souls interpenetrate. So even if an animal does not have an independent intellectual soul, the two lower souls do have an intellectual component. That is the manifestation of what we called an “angelic intelligence”.

Lewis is just speculating. Any “self” he seems to observe does not belong to the animal as an individual begin, but may perhaps be a vestige of that angelic intelligence. Of course animals suffer since they are sentient beings. We're not in agreement with Descartes on that point.


\end{sffamily}\end{footnotesize}
